\chapter{CPU Scheduling}

\section{Chapter Overview}

CPU Scheduling is one of the most important Operating Systems topics for Software Development Engineer interviews. It combines theoretical concepts with numerical problem-solving and is frequently asked by Amazon, Google, Microsoft, Meta, Uber, Atlassian, Apple, Bloomberg, and many other technology companies.

This chapter covers scheduling fundamentals, scheduling metrics, major scheduling algorithms, numerical techniques, interview tricks, and quick-solving strategies.

\section{Learning Objectives}

After completing this chapter, you should be able to:

\begin{itemize}
\item Understand why CPU scheduling is necessary.
\item Calculate scheduling performance metrics.
\item Differentiate preemptive and non-preemptive scheduling.
\item Solve scheduling numericals efficiently.
\item Analyze FCFS, SJF, SRTF, Priority, Round Robin, MLQ, and MLFQ.
\item Explain starvation and aging.
\item Compare scheduling algorithms.
\item Answer interview questions confidently.
\end{itemize}

%====================================================
\section{CPU Scheduling Basics}
%====================================================

\subsection{Definition}

CPU Scheduling is the process of selecting a process from the ready queue and allocating the CPU to it.

\subsection{Why Scheduling is Needed}

At any instant:

\begin{itemize}
\item Multiple processes may be ready.
\item Only one process can use a CPU core.
\item OS must decide who gets CPU access.
\end{itemize}

\subsection{Scheduling Cycle}

\begin{center}
Ready Queue $\rightarrow$ Scheduler $\rightarrow$ CPU Execution
\end{center}

\begin{figure}[H]
\centering
\begin{tikzpicture}

\node[draw] (rq) {Ready Queue};
\node[draw,right=4cm of rq] (cpu) {CPU};

\draw[->,thick] (rq) -- node[above]{Scheduler} (cpu);

\end{tikzpicture}
\caption{Basic CPU Scheduling}
\end{figure}

%====================================================
\section{Scheduling Goals}
%====================================================

A good scheduler should:

\begin{itemize}
\item Maximize CPU utilization.
\item Maximize throughput.
\item Minimize waiting time.
\item Minimize turnaround time.
\item Minimize response time.
\item Ensure fairness.
\item Avoid starvation.
\end{itemize}

%====================================================
\section{Throughput}
%====================================================

\subsection{Definition}

Number of processes completed per unit time.

\[
Throughput =
\frac{\text{Number of completed processes}}
{\text{Time}}
\]

\subsection{Example}

50 processes completed in 10 seconds.

\[
Throughput = \frac{50}{10}=5
\]

processes/second.

%====================================================
\section{Turnaround Time}
%====================================================

\subsection{Definition}

Total time spent by a process in the system.

\[
TAT = Completion - Arrival
\]

\subsection{Example}

Arrival = 2

Completion = 15

\[
TAT=15-2=13
\]

%====================================================
\section{Waiting Time}
%====================================================

\subsection{Definition}

Time spent waiting in the ready queue.

\[
WT=TAT-Burst
\]

\subsection{Example}

TAT = 20

Burst = 8

\[
WT=20-8=12
\]

%====================================================
\section{Response Time}
%====================================================

\subsection{Definition}

Time from arrival until first CPU allocation.

\[
RT = First\ Start - Arrival
\]

\subsection{Importance}

Critical for interactive systems.

%====================================================
\section{CPU Utilization}
%====================================================

\subsection{Definition}

Percentage of time CPU remains busy.

\[
CPU\ Utilization=
\frac{CPU\ Busy\ Time}
{Total\ Time}
\times100
\]

\subsection{Goal}

Keep CPU utilization as high as possible.

%====================================================
\section{Preemptive Scheduling}
%====================================================

\subsection{Definition}

OS can forcibly take CPU away from a running process.

\subsection{Examples}

\begin{itemize}
\item Round Robin
\item SRTF
\item Preemptive Priority
\end{itemize}

\subsection{Advantages}

\begin{itemize}
\item Better responsiveness
\item Fair CPU allocation
\end{itemize}

\subsection{Disadvantages}

\begin{itemize}
\item More context switches
\item Higher overhead
\end{itemize}

%====================================================
\section{Non-Preemptive Scheduling}
%====================================================

\subsection{Definition}

Once CPU is allocated, process keeps CPU until completion or blocking.

\subsection{Examples}

\begin{itemize}
\item FCFS
\item SJF
\item Non-preemptive Priority
\end{itemize}

\subsection{Advantages}

\begin{itemize}
\item Simpler
\item Less overhead
\end{itemize}

\subsection{Disadvantages}

\begin{itemize}
\item Poor responsiveness
\item Convoy effect
\end{itemize}

%====================================================
\section{First Come First Serve (FCFS)}
%====================================================

\subsection{Definition}

Processes are executed in order of arrival.

\subsection{Working}

Queue behaves like FIFO.

\subsection{Advantages}

\begin{itemize}
\item Simple
\item Easy implementation
\end{itemize}

\subsection{Disadvantages}

\begin{itemize}
\item Convoy effect
\item Poor average waiting time
\end{itemize}

\subsection{Complexity}

\[
O(1)
\]

\subsection{Example}

\begin{table}[H]
\centering
\begin{tabular}{|c|c|}
\hline
Process & Burst \\
\hline
P1 & 5\\
P2 & 3\\
P3 & 2\\
\hline
\end{tabular}
\end{table}

Gantt Chart:

\begin{center}
0 \quad P1 \quad 5 \quad P2 \quad 8 \quad P3 \quad 10
\end{center}

Waiting Times:

\[
P1=0
\]

\[
P2=5
\]

\[
P3=8
\]

Average WT:

\[
\frac{0+5+8}{3}=4.33
\]

%====================================================
\section{Shortest Job First (SJF)}
%====================================================

\subsection{Definition}

Shortest burst process executes first.

\subsection{Working}

Select minimum burst time.

\subsection{Advantages}

\begin{itemize}
\item Minimum average waiting time
\end{itemize}

\subsection{Disadvantages}

\begin{itemize}
\item Burst prediction difficult
\item Starvation possible
\end{itemize}

\subsection{Complexity}

\[
O(n)
\]

\subsection{Example}

Processes:

\[
P1=6,\;
P2=2,\;
P3=8,\;
P4=3
\]

Order:

\[
P2 \rightarrow P4 \rightarrow P1 \rightarrow P3
\]

Gantt Chart:

\begin{center}
0 2 5 11 19
\end{center}

Average WT:

\[
\frac{0+2+5+11}{4}
=4.5
\]

%====================================================
\section{Shortest Remaining Time First (SRTF)}
%====================================================

\subsection{Definition}

Preemptive version of SJF.

\subsection{Working}

Always execute process with smallest remaining burst.

\subsection{Advantages}

\begin{itemize}
\item Optimal average waiting time
\end{itemize}

\subsection{Disadvantages}

\begin{itemize}
\item High context switching
\item Starvation possible
\end{itemize}

\subsection{Complexity}

\[
O(n)
\]

\subsection{Example}

Arrival:

\[
P1=(0,8)
\]

\[
P2=(1,4)
\]

At time 1:

P2 preempts P1.

Gantt:

\begin{center}
P1 | P2 | P1
\end{center}

%====================================================
\section{Priority Scheduling}
%====================================================

\subsection{Definition}

Highest priority process executes first.

\subsection{Working}

Priority determines scheduling order.

\subsection{Advantages}

\begin{itemize}
\item Important tasks run first
\end{itemize}

\subsection{Disadvantages}

\begin{itemize}
\item Starvation
\end{itemize}

\subsection{Complexity}

\[
O(n)
\]

\subsection{Example}

\begin{table}[H]
\centering
\begin{tabular}{|c|c|}
\hline
Process & Priority \\
\hline
P1 & 3\\
P2 & 1\\
P3 & 2\\
\hline
\end{tabular}
\end{table}

Order:

\[
P2 \rightarrow P3 \rightarrow P1
\]

%====================================================
\section{Round Robin Scheduling}
%====================================================

\subsection{Definition}

Processes receive CPU for a fixed time quantum.

\subsection{Working}

After quantum expires:

\begin{itemize}
\item Process preempted
\item Moved to queue end
\end{itemize}

\subsection{Advantages}

\begin{itemize}
\item Fairness
\item Good response time
\end{itemize}

\subsection{Disadvantages}

\begin{itemize}
\item High context switching
\item Quantum selection critical
\end{itemize}

\subsection{Complexity}

\[
O(1)
\]

\subsection{Example}

Quantum = 2

\[
P1=5
\]

\[
P2=4
\]

Gantt:

\begin{center}
P1 P2 P1 P2 P1
\end{center}

%====================================================
\section{Multilevel Queue Scheduling}
%====================================================

\subsection{Definition}

Ready queue divided into multiple fixed queues.

Example:

\begin{itemize}
\item System Processes
\item Interactive Processes
\item Batch Processes
\end{itemize}

\subsection{Working}

Each queue has its own scheduler.

\subsection{Advantages}

\begin{itemize}
\item Separation of workloads
\end{itemize}

\subsection{Disadvantages}

\begin{itemize}
\item Inflexible
\item Starvation possible
\end{itemize}

\subsection{Complexity}

Depends on queue structure.

%====================================================
\section{Multilevel Feedback Queue (MLFQ)}
%====================================================

\subsection{Definition}

Processes can move between queues.

\subsection{Working}

\begin{itemize}
\item New jobs start at top queue.
\item CPU-intensive jobs move downward.
\item Interactive jobs remain near top.
\end{itemize}

\subsection{Advantages}

\begin{itemize}
\item Adaptive
\item Excellent responsiveness
\end{itemize}

\subsection{Disadvantages}

\begin{itemize}
\item Complex implementation
\end{itemize}

\subsection{Complexity}

Implementation dependent.

\subsection{Example}

\[
Q1: RR(4)
\]

\[
Q2: RR(8)
\]

\[
Q3: FCFS
\]

%====================================================
\section{Dispatcher}
%====================================================

\subsection{Definition}

Dispatcher transfers CPU control to selected process.

Responsibilities:

\begin{itemize}
\item Context switching
\item User mode transition
\item Jump to process location
\end{itemize}

\subsection{Dispatch Latency}

Time taken by dispatcher to stop one process and start another.

%====================================================
\section{Context Switch Impact}
%====================================================

Context switching introduces overhead.

\[
Useful\ Work = Total\ Time - Context\ Switch\ Time
\]

Excessive context switching reduces:

\begin{itemize}
\item Throughput
\item CPU efficiency
\end{itemize}

%====================================================
\section{Starvation}
%====================================================

\subsection{Definition}

A process waits indefinitely because other processes continuously receive CPU.

\subsection{Common In}

\begin{itemize}
\item Priority Scheduling
\item SJF
\item SRTF
\end{itemize}

%====================================================
\section{Aging}
%====================================================

\subsection{Definition}

Gradually increasing process priority over time.

\subsection{Purpose}

Prevents starvation.

\subsection{Example}

Priority increases every 5 seconds.

%====================================================
\section{Comparative Summary}
%====================================================

\begin{table}[H]
\centering
\scriptsize
\begin{tabular}{|p{3cm}|p{2cm}|p{2cm}|p{3cm}|}
\hline
Algorithm & Preemptive & Starvation & Response Time \\
\hline
FCFS & No & No & Poor \\
\hline
SJF & No & Yes & Good \\
\hline
SRTF & Yes & Yes & Excellent \\
\hline
Priority & Both & Yes & Good \\
\hline
RR & Yes & No & Excellent \\
\hline
MLQ & Usually & Possible & Good \\
\hline
MLFQ & Yes & Rare & Excellent \\
\hline
\end{tabular}
\caption{Scheduling Algorithm Comparison}
\end{table}

%====================================================
\section{Common Interview Traps}
%====================================================

\begin{enumerate}
\item Forgetting arrival times.
\item Confusing turnaround and waiting time.
\item Using wrong quantum in RR.
\item Ignoring preemption in SRTF.
\item Incorrect completion times.
\item Missing starvation discussion.
\item Not mentioning aging.
\item Forgetting context switch overhead.
\end{enumerate}

%====================================================
\section{How to Solve Scheduling Numericals Quickly}
%====================================================

\subsection{Step 1}

Create process table.

\subsection{Step 2}

Draw Gantt chart.

\subsection{Step 3}

Find completion times.

\subsection{Step 4}

Calculate:

\[
TAT = CT - AT
\]

\[
WT = TAT - BT
\]

\[
RT = FirstStart - AT
\]

\subsection{Step 5}

Compute averages.

\subsection{Golden Interview Formula}

\[
WT = TAT - BT
\]

This formula solves most scheduling problems quickly.

%====================================================
\section{Key Takeaways}
%====================================================

\begin{itemize}
\item Scheduling selects the next CPU process.
\item FCFS is simple but suffers convoy effect.
\item SJF minimizes average waiting time.
\item SRTF is preemptive SJF.
\item Priority scheduling can starve low-priority jobs.
\item Aging prevents starvation.
\item Round Robin provides fairness.
\item MLFQ is widely used in modern systems.
\item Dispatcher performs context switches.
\item Numerical problems require accurate Gantt charts.
\end{itemize}

%====================================================
\section{25 Interview Questions with Answers}
%====================================================

\begin{enumerate}
\item What is CPU scheduling?\\
Answer: Selecting a process for CPU execution.

\item What is throughput?\\
Answer: Completed processes per unit time.

\item Define turnaround time.\\
Answer: Completion time minus arrival time.

\item Define waiting time.\\
Answer: Time spent in ready queue.

\item Define response time.\\
Answer: First CPU allocation delay.

\item What is CPU utilization?\\
Answer: Percentage CPU remains busy.

\item FCFS stands for?\\
Answer: First Come First Serve.

\item Major issue in FCFS?\\
Answer: Convoy effect.

\item Which algorithm minimizes average waiting time?\\
Answer: SJF.

\item What is SRTF?\\
Answer: Preemptive SJF.

\item What causes starvation?\\
Answer: Continuous postponement.

\item How is starvation prevented?\\
Answer: Aging.

\item What is aging?\\
Answer: Gradual priority increase.

\item Round Robin is?\\
Answer: Preemptive.

\item What is time quantum?\\
Answer: Fixed CPU slice.

\item Very small quantum causes?\\
Answer: Excessive context switching.

\item Very large quantum becomes?\\
Answer: FCFS.

\item Priority scheduling problem?\\
Answer: Starvation.

\item What is dispatcher?\\
Answer: CPU transfer mechanism.

\item Dispatch latency?\\
Answer: Dispatcher overhead time.

\item Which algorithm is best for interactive systems?\\
Answer: Round Robin.

\item Which algorithm adapts dynamically?\\
Answer: MLFQ.

\item Which algorithm uses multiple queues?\\
Answer: MLQ.

\item What is context switching?\\
Answer: Process state change.

\item Why is scheduling important?\\
Answer: Efficient CPU utilization.
\end{enumerate}

%====================================================
\section{25 MCQs with Answers}
%====================================================

\begin{enumerate}
\item FCFS is:
A) Preemptive
B) Non-preemptive
Answer: B

\item SRTF is:
Answer: Preemptive SJF

\item Turnaround Time =
Answer: Completion - Arrival

\item Waiting Time =
Answer: TAT - Burst

\item Response Time measures:
Answer: First CPU allocation delay

\item RR uses:
Answer: Time Quantum

\item Starvation occurs in:
Answer: Priority Scheduling

\item Aging prevents:
Answer: Starvation

\item Dispatcher performs:
Answer: Context Switch

\item CPU utilization should be:
Answer: High

\item Throughput should be:
Answer: High

\item Waiting time should be:
Answer: Low

\item Round Robin is:
Answer: Preemptive

\item Convoy effect occurs in:
Answer: FCFS

\item SJF gives:
Answer: Minimum Average Waiting Time

\item Interactive systems prefer:
Answer: RR

\item MLFQ stands for:
Answer: Multilevel Feedback Queue

\item Dispatch latency is:
Answer: Dispatcher Time

\item Priority scheduling can cause:
Answer: Starvation

\item MLFQ allows:
Answer: Queue Migration

\item Small quantum causes:
Answer: High Overhead

\item Large quantum resembles:
Answer: FCFS

\item Scheduler selects:
Answer: Next Process

\item Throughput unit:
Answer: Processes/Time

\item SRTF chooses:
Answer: Shortest Remaining Burst
\end{enumerate}

%====================================================
\section{Revision Notes}
%====================================================

\begin{itemize}
\item CPU Scheduling selects next process.
\item Throughput should increase.
\item Waiting time should decrease.
\item TAT = CT - AT.
\item WT = TAT - BT.
\item RT = First Start - AT.
\item FCFS suffers convoy effect.
\item SJF minimizes average waiting time.
\item SRTF is preemptive SJF.
\item RR uses time quantum.
\item Aging prevents starvation.
\item MLFQ is adaptive.
\end{itemize}

%====================================================
\section{Cheat Sheet}
%====================================================

\begin{center}
\begin{tabular}{|p{5cm}|p{8cm}|}
\hline
TAT & CT - AT \\
\hline
WT & TAT - BT \\
\hline
RT & First Start - AT \\
\hline
FCFS & FIFO Scheduling \\
\hline
SJF & Smallest Burst First \\
\hline
SRTF & Smallest Remaining Burst First \\
\hline
Priority & Highest Priority First \\
\hline
RR & Time Quantum Scheduling \\
\hline
MLQ & Fixed Multiple Queues \\
\hline
MLFQ & Dynamic Multiple Queues \\
\hline
Starvation & Indefinite Waiting \\
\hline
Aging & Priority Increase \\
\hline
Dispatcher & CPU Transfer Unit \\
\hline
Context Switch & Save + Restore State \\
\hline
\end{tabular}
\end{center}